\chapter{Heel Pain}

\section*{Definition}
Pain is localized to the posterior or inferior aspect of the heel, commonly caused by mechanical stress on the plantar fascia or Achilles tendon insertion.

\section*{What Happens in Your Body}
Plantar fasciitis results from repetitive microtrauma and degeneration of the plantar fascia origin at the medial calcaneal tuberosity. Achilles tendinopathy involves degeneration of the Achilles tendon fibers from overuse, poor footwear, or biomechanical abnormalities. Heel pain is characteristically worse with the first steps in the morning or after prolonged sitting.

\section*{Causes}
\begin{itemize}
  \item Plantar fasciitis (most common cause)
  \item Achilles tendinopathy or insertional tendinitis
  \item Calcaneal stress fracture
  \item Retrocalcaneal bursitis
  \item Heel spur syndrome
  \item Sever disease (calcaneal apophysitis in children)
  \item Fat pad atrophy (especially in elderly)
  \item Tarsal tunnel syndrome
  \item Gout or pseudogout
  \item Seronegative spondyloarthropathy (ankylosing spondylitis, reactive arthritis)
\end{itemize}

\section*{When to See a Doctor}
See your doctor if:
\begin{itemize}
  \item Symptoms are severe or getting worse over time
  \item Severe heel pain with inability to bear weight, fever, redness, or swelling over the heel
  \item You have other concerning symptoms such as fever, weight loss, or bleeding
\end{itemize}

\section*{Related Symptoms}
\begin{itemize}
  \item Arch pain
  \item Achilles tendon pain
  \item Ankle pain
  \item Foot numbness
  \item Morning stiffness
\end{itemize}
\textbf{Important:} If this symptom persists, your doctor will check for serious underlying causes.

\section*{Self-Care / Home Management}
\begin{itemize}
  \item Rest and avoid activities that make it worse.
  \item Maintain adequate hydration and a balanced diet.
  \item Over-the-counter remedies may provide symptomatic relief (consult a pharmacist or doctor).
  \item Monitor symptoms and seek professional advice if they persist or worsen.
  \item Avoid self-medicating with prescription medications without medical guidance.
\end{itemize}

\section*{How Your Doctor Will Evaluate You}
\begin{itemize}
  \item A thorough history and physical examination are the first steps in evaluation.
  \item Laboratory tests (blood work, urinalysis) may be ordered based on clinical suspicion.
  \item Imaging studies (X-ray, ultrasound, CT, MRI) may be indicated when structural pathology is suspected.
  \item Specialist referral is arranged when the diagnosis remains uncertain or requires specific expertise.
\end{itemize}

\section*{Treatment Options}
\begin{itemize}
  \item Treatment targets the underlying cause identified through diagnostic evaluation.
  \item Supportive care and symptom management are provided while awaiting definitive therapy.
  \item Medications, lifestyle modifications, and physical therapies may be prescribed as appropriate.
  \item Follow-up ensures treatment effectiveness and allows timely adjustments to the care plan.
\end{itemize}

\clearpage